\section{Coupled cube barriers with optimal parameter}
\label{sec:coupled-barriers}
The barriers of \cref{sec:barrier-dependence} are separable. We show
that the worst-case ratios, between endpoints and at equal accuracy,
stay at least $\Gamma_r$ for some coupling terms that make the Hessian
non-diagonal: we prove this for convex couplings bounded near one
facet, give two coupled barriers with the optimal parameter, and prove a
matching upper bound for one of them. Throughout, $K=(-1,1)^r$ and
$U(x)=\sum_i b(x_i)$. Every standard self-concordant barrier on $K$ has
parameter at least $r$; this classical bound follows because the
vertices of $K$ are simple \cite[Proposition 2.3.6]{NesterovNemirovskii1994}.

\subsection{A lower bound when the coupling is bounded near a facet}
For $\delta\in(0,1)$, the facet collar of width $\delta$ is
$\{x\in K:x_1\geq1-\delta\}$.
\begin{theorem}[Lower bound under bounded coupling near a facet]
\label{thm:facet-coupling}
Let $F=U+G$ be a standard self-concordant barrier, with $G\in C^2(K)$
convex. Suppose that, after a signed coordinate permutation, some
$\delta\in(0,1)$ and finite $M,C$ satisfy
\begin{equation}
 \norm{\nabla G(x)}_\infty\leq M,\quad
 0\preceq G''(x)\preceq CI
 \quad\text{for every }x\in K\text{ with }x_1\geq1-\delta.
 \label{eq:facet-collar}
\end{equation}
Let $x_*$ be the analytic center of $F$. Then
\begin{equation}
 \sup_{w>0,\eta>0}
 \frac{L_{\CP,F}(0,\eta)}{d_F(x_*,x_w(\eta))}\geq\Gamma_r,
 \qquad
 \sup_{w>0,\,0<\eps<g_w(x_*)}
 \frac{L_{\CP,F}(\eps)}{L_{\opt,F}(\eps)}\geq\Gamma_r.
 \label{eq:facet-tax}
\end{equation}
In the second ratio, distances to accuracy are measured from $x_*$.
The weights are positive in the permuted coordinates. We do not claim a matching upper bound
for all $G$ satisfying \eqref{eq:facet-collar}.
\end{theorem}
\begin{proof}
The bounded cube and barrier blow-up give an interior minimizer,
unique because $F''\succeq U''\succ0$. Let
$d_i=\sqrt i-\sqrt{i-1}$, choose $w_i=e^{Md_i}$ with
$M\to\infty$, and write $x_M(s)$ for the center at $\eta=e^s$,
terminating at $s=0$. In this proof only, denote the gradient bound
in \eqref{eq:facet-collar} by $M_G$ to distinguish it from $M$.
Set $z_i=e^{s+Md_i}$. For each fixed $x_{-1}$, convexity implies
\[
 x_1\leq1-\delta\quad\Longrightarrow\quad
 \partial_1G(x_1,x_{-1})\leq
 \partial_1G(1-\delta,x_{-1})\leq M_G.
\]
Choose $Z>M_G+b'(1-\delta)$. The equation
$b'(x_1)+\partial_1G=z_1$ shows that $z_1\geq Z$ forces entry
into the collar. As $z_1$ increases, the path stays there.

In the collar, put $D=U''(x_M)$ and $H=F''(x_M)$.
Differentiated stationarity gives $H\dot x_M=z$ and speed squared
$z^TH^{-1}z$. For any coordinate subset $A$, the variational
formula for this inverse quadratic form, restricted to vectors
supported on $A$, gives
\begin{equation}
 z^TH^{-1}z\geq z_A^TH_{AA}^{-1}z_A
 \geq\sum_{i\in A}\frac{z_i^2}{b''(x_{M,i})+C}.
 \label{eq:coupled-prefix-speed}
\end{equation}
Also $b'(x_{M,i})=z_i-\partial_iG$ with
$|\partial_iG|\leq M_G$. Because
$(b')^2/b''=2x^2/(1+x^2)\to1$ as $x\uparrow1$, increasing
$Z$ makes every summand in \eqref{eq:coupled-prefix-speed}
at least $1-\zeta$ whenever $z_i\geq Z$, for any fixed
$\zeta>0$. The activation times
$\tau_i=\log Z-Md_i$ are strictly ordered, and all lie below zero
for large $M$. Integrating over their successive prefix intervals gives
\begin{equation}
 L_{\CP,F}(0,1)\geq
 \sqrt{1-\zeta}\,[M\Gamma_r^2-\sqrt r\log Z].
 \label{eq:facet-arc-lower}
\end{equation}
At termination, stationarity and the bounded gradient imply
$y_i:=\rho(x_{M,i}(0))=Md_i+O_{r,M_G}(1)$, by
\eqref{eq:profile-tails}; hence $\norm y=M\Gamma_r+O(1)$.

Fix an interior point $p$ with $p_1=1-\delta/2$ and a fixed
compact interior curve from $x_*$ to $p$. Its $F$-length is finite.
The straight segment from $\rho(p)$ to $y$ in product metric
coordinates has monotone physical coordinates. For large $M$ its
first coordinate stays above $1-\delta$, so it lies in the collar.
Its $U$-length is $M\Gamma_r+O(1)$ and its Euclidean total
variation is at most $2r$. Its extra $G$-length is therefore at most
$2r\sqrt C$. Using $\sqrt{a+b}\leq\sqrt a+\sqrt b$ gives
\begin{equation}
 d_F(x_*,x_M(0))\leq M\Gamma_r+O_{F,r,M_G,C}(1).
 \label{eq:facet-distance-upper}
\end{equation}
Divide \eqref{eq:facet-arc-lower} by this bound and let first
$M\to\infty$, then $\zeta\downarrow0$. This proves the first
assertion, with analytic-center initialization interpreted as the
limit $s\to-\infty$. Near that limit the path has finite length
because $x(\eta)=x_*+\eta F''(x_*)^{-1}w+o(\eta)$.

For each such objective and terminal center, choose
$\eps=g_w(x_M(0))>0$. The identity
\[
 \frac{\dd}{\dd\eta}w^Tx_w(\eta)=w^TF''(x_w)^{-1}w>0
\]
shows this is the first accurate center and
$\eps<g_w(x_*)$. Since
$L_{\opt,F}(\eps)\leq d_F(x_*,x_M(0))$, the same construction
proves the second assertion.
\end{proof}
The proof needs the bounds in \eqref{eq:facet-collar} on the whole
collar, not only near a vertex: it uses them after the first coordinate
enters the collar, while the others are still far from their facets.
Hence, for convex $G$, a worst-case ratio bounded independently of $r$
requires \eqref{eq:facet-collar} to fail for every signed permutation
and every $\delta$, that is, derivatives of $G$ unbounded in every
facet collar. This is a necessary condition; we do
not claim that such a $G$ exists.

\subsection{Two explicit coupled barriers with the optimal parameter}
\begin{theorem}[Dense and radial coupling at the cube parameter]
\label{thm:optimal-couplings}
For $r\geq2$, each of the barriers
\begin{align}
 F_{\mathrm d}(x)&=U(x)+\frac1{8r}(\mathbf1^Tx)^2,
       \label{eq:dense-barrier}\\
 F_{\lambda,c}(x)&=U(x)-\lambda\log(c-\norm x^2),
 \qquad\lambda\geq1,\quad c-r\geq4\lambda,
       \label{eq:radial-barrier}
\end{align}
is standard self-concordant with exact parameter $r$ and has
nonzero mixed Hessian entries on an open subset of the cube.
The dense barrier is invariant under permutations and simultaneous
sign reversal. The radial family is invariant under every signed
permutation. For each fixed choice of their parameters, both suprema
in \eqref{eq:facet-tax} are at least $\Gamma_r$.
\end{theorem}
\begin{proof}
Write $p=\nabla U$ and $D=U''$. For the dense example put
$\beta=1/(4r)$, $t=\mathbf1^Tx$,
$A=\mathbf1^TD^{-1}\mathbf1$, $B=\mathbf1^TD^{-1}p$, and
$Q=p^TD^{-1}p$. Sherman--Morrison gives
\[
 \norm{\nabla F_{\mathrm d}}_{F_{\mathrm d},x,*}^2
 =Q+\frac{\beta(2tB-B^2)+\beta^2t^2A}{1+\beta A}.
\]
Let $m=r-Q=\sum_i(1-x_i^2)/(1+x_i^2)>0$.
The explicit interval derivatives give $|B|\leq m$, $A\leq m/2$,
and $|t|<r$. Dropping $-\beta B^2$ bounds the correction by
$(2\beta r+\beta^2r^2/2)m=17m/32<m$.
Thus the parameter is at most $r$. The quadratic has zero third
derivative and positive semidefinite Hessian, so it preserves
standard self-concordance; $U$ supplies boundary blow-up.

For the radial case put $S=c-\norm x^2$ and $a=2\lambda/S$.
Then $0<a<1/2$ and
\begin{equation}
 \nabla F_{\lambda,c}=p+ax,\qquad
 F_{\lambda,c}''=D+aI+(a^2/\lambda)xx^T.
 \label{eq:radial-derivatives}
\end{equation}
The ball term is an affine rescaling of the rank-one spectral-ball
barrier in \cref{prop:standard-parameters}; multiplying by
$\lambda\geq1$, restricting to the cube, and summing preserve
standard self-concordance. Inverse Loewner order gives
\[
 \norm{\nabla F_{\lambda,c}}_{F_{\lambda,c},x,*}^2
 \leq\sum_i\frac{(p_i+ax_i)^2}{D_i+a}\leq r.
\]
For the last inequality set $y=x_i^2$. After multiplying by
$(1-y)^2$, the difference between denominator and numerator
in its $i$th summand is
\[
 (1-y)\,[2+a(1-5y)-a^2y(1-y)].
\]
The bracket decreases on $[0,1]$, since its derivative is
$a(2ay-a-5)<0$, and its value at one is $2-4a\geq0$.

For both examples, along $x=t\mathbf1$ as $t\uparrow1$ the
added gradient and Hessian stay bounded, while the interval gradient
and Hessian diverge. The ratio of the squared first to the second directional derivative
along $\mathbf1$ tends to $r$; explicitly it is
$r(p_0+at)^2/[D_0+a+(a^2/\lambda)rt^2]$ for the radial example.
Thus both exact parameters equal $r$.
The stated symmetries follow from the formulas. The mixed Hessians
are $\beta$ in the dense case and
$(a^2/\lambda)x_ix_j$ in the radial case.
Finally the dense coupling obeys
$\norm{\nabla G}_\infty\leq1/4$, $G''\preceq I/4$.
The radial coupling obeys
\[
 \norm{\nabla G}_\infty<1/2,\qquad
 0\preceq G''\preceq(1/2+r/(4\lambda))I.
\]
These global bounds imply \eqref{eq:facet-collar}; apply
\cref{thm:facet-coupling}.
\end{proof}
Adding the parameters of the two terms in \eqref{eq:radial-barrier}
gives only $r+\lambda$, whereas the exact parameter is $r$; with
$c=r+4\lambda$ and large $\lambda$ the gap is arbitrarily large. This
requires $c$ to vary, since for fixed $r$ and $c$ our sufficient
condition allows only $\lambda\leq(c-r)/4$.
Quadratic regularization that preserves the parameter of the box
barrier is known: Castro and Cuesta~\cite[Section 4.1]{CastroCuesta2011}
prove this for suitably bounded diagonal regularizations. We use our two examples only for the central-path comparisons and do
not claim that parameter-preserving regularization is new.

\subsection{A uniform matching upper bound for the radial family}
\begin{theorem}[Upper bound for the radial family]
\label{thm:radial-upper}
For every barrier \eqref{eq:radial-barrier} and every nonzero linear
objective with $k$ nonzero coordinates, every central arc satisfies
\begin{equation}
 L_{\CP,F}(\eta_0,\eta_1)\leq
 B_{\mathrm r}\Gamma_k d_F(x(\eta_0),x(\eta_1)),\qquad
 B_{\mathrm r}=\sqrt{\frac{11}{8}}\frac{25}{14}.
 \label{eq:radial-endpoint-upper}
\end{equation}
For the central path from the analytic center $0$ and
$0<\eps<\sum_i|w_i|$,
\begin{equation}
 L_{\CP,F}(\eps)\leq
 \frac54 B_{\mathrm r}c_\star\Gamma_k L_{\opt,F}(\eps).
 \label{eq:radial-accuracy-upper}
\end{equation}
The constants do not depend on $r$, $\lambda$, or $c$; we do not
claim that they are sharp. Together with \cref{thm:facet-coupling},
these inequalities show that, for every barrier in this family, both
worst-case ratios over positive objectives are $\Theta(\sqrt{\log r})$.
\end{theorem}
\begin{proof}
Let $t=\sum_i(1-x_i^2)$, so $S\geq4\lambda+t$.
The interval Hessian satisfies $D\succeq2I$ and
$x^TD^{-1}x\leq t/2$. The two coupling terms in
\eqref{eq:radial-derivatives} satisfy
\[
 aI\preceq\tfrac14D,\qquad
 \frac{a^2}{\lambda}xx^T\preceq
 \frac{2\lambda t}{(4\lambda+t)^2}D\preceq\tfrac18D.
\]
The last inequality follows from $(t-4\lambda)^2\geq0$.
Consequently, throughout the whole cube,
\begin{equation}
 U''\preceq F''\preceq\frac{11}{8}U''.
 \label{eq:radial-metric-comparison}
\end{equation}
In particular $d_F\geq d_U$, even for paths that leave the subspace of
active (nonzero-weight) coordinates.

By signed-permutation invariance take the nonzero weights positive
and decreasing. The stationary equations at $\eta=e^s$ are
\begin{equation}
 z_i=e^sw_i=x_i(b_i+a),\qquad b_i=2/(1-x_i^2).
 \label{eq:radial-central-equations}
\end{equation}
Zero-weight coordinates are zero. Every other coordinate is positive,
and their order is that of $w_i$, since $x\mapsto x(b(x)+a)$
is strictly increasing for fixed $a$.
Differentiating with respect to $s$ and solving for $a'$ gives
\begin{align}
 (D_i+a)x_i'+x_i a'&=z_i,\notag\\
 a'&=\frac{(a^2/\lambda)K}{1+(a^2/\lambda)L},\qquad
 K=\sum_i\frac{x_i^2(b_i+a)}{D_i+a},\quad
 L=\sum_i\frac{x_i^2}{D_i+a}.
 \label{eq:radial-a-prime}
\end{align}
For $y=x_i^2$, the inequality
$y(b_i+a)/(D_i+a)\leq1-y$ reduces, after multiplication
by $1-y$, to $2+a(1-2y)(1-y)>0$.
Thus $K\leq t$ and
\begin{equation}
 0\leq a'\leq\frac{4\lambda t}{(4\lambda+t)^2}\leq\frac14.
 \label{eq:radial-a-speed}
\end{equation}
Put $q_i=b'(x_i)=b_ix_i$ on active coordinates and define
\[
 \theta_i=(\log q_i)'=
 \frac{D_i(b_i+a-a')}{(D_i+a)b_i},\qquad
 v_i=\frac{q_i}{\sqrt{D_i}}=
 \frac{\sqrt2x_i}{\sqrt{1+x_i^2}}.
\]
Since $D_i,b_i\geq2$, $a\leq1/2$, and $a'\leq1/4$,
\begin{equation}
 \frac7{10}=\frac45\frac78\leq\theta_i\leq\frac54,
 \qquad y_i':=(\rho(x_i))'=\theta_i v_i.
 \label{eq:radial-ordered-speed}
\end{equation}
In particular all coordinates increase. The $v_i$ are nonnegative
and decreasing in index, since the $x_i$ are ordered.
Let $d_i=\sqrt i-\sqrt{i-1}$. Combining the metric comparison,
the bounds on $\theta_i$, and the prefix inequality yields
\begin{align*}
 L_{\CP,F}&\leq\sqrt{11/8}\int\norm{y'}_2\dd s
 \leq\sqrt{11/8}\frac54\int\norm{v}_2\dd s\\
 &\leq\sqrt{11/8}\frac54\sum_{i=1}^k d_i\int v_i\dd s
 \leq B_{\mathrm r}\sum_{i=1}^k d_i\Delta y_i
 \leq B_{\mathrm r}\Gamma_k\norm{\Delta y}_2.
\end{align*}
The last norm equals the full-cube $U$-distance of the two centers
and is at most their $F$-distance. This proves
\eqref{eq:radial-endpoint-upper}, including arcs that start at the
analytic center.

For accuracy, compare with the standard $U$ center for the same
weights. Equation \eqref{eq:radial-central-equations} gives
\begin{equation}
 \frac45z_i\leq q_i\leq z_i.
 \label{eq:radial-q-comparison}
\end{equation}
Let $\eta_U$ be the first accurate $U$ parameter. The radial
center at $\bar\eta=5\eta_U/4$ is coordinatewise at least this
$U$ center and is therefore accurate. By coordinate monotonicity the first
accurate radial center occurs no later.
For the function $H$ in \eqref{eq:standard-profile},
$v=H'$ has $v'\leq v$ because $b'''\geq0$ on $(0,1)$.
Integration from $-\infty$ gives $v\leq H$, and hence
$H(s+\log a)\leq aH(s)$ for $a\geq1$.
The upper bound in \eqref{eq:radial-q-comparison} consequently
shows that the radial $\rho$ coordinates at $\bar\eta$ are at
most $5/4$ times the coordinates of the first accurate $U$ center.
The proof of \cref{thm:scalar-dilation} bounds the latter vector's
norm by $c_\star L_{\opt,U}(\eps)$.
Finally $L_{\opt,U}(\eps)\leq L_{\opt,F}(\eps)$ follows from
global metric domination. Apply the direct bound
$L_{\CP,F}(0,\bar\eta)\leq B_{\mathrm r}\Gamma_k\norm{y(\bar\eta)}$
proved above, obtaining \eqref{eq:radial-accuracy-upper}.
\end{proof}

\subsection{The dyadic separation for the radial family}
\begin{corollary}[Uniform radial discrete comparison]
\label{cor:radial-discrete}
Use the same dyadic objective and tolerance \eqref{eq:dyadic-family}
with any $F_{\lambda,c}$ satisfying \eqref{eq:radial-barrier};
the parameters may vary with $r$. Uniformly over this family,
\[
 L_{\opt,F}(\eps_r)=\Theta(r\sqrt{\log r}),\qquad
 L_{\CP,F}(\eps_r)=\Theta(r\log r).
\]
Let $0<R<1$ and $\delta_r=o(r^{2/3})$. Every forward-$R$ sequence
with arbitrary labels that starts at zero, ends at a point with gap at
most $\eps_r$, and stays within $F$-distance $\delta_r$ of its labeled
central points satisfies \eqref{eq:growing-tube}. An explicit
unrestricted sequence reaches accuracy $\eps_r$ in
$O_R(r\sqrt{\log r})$ steps. In particular, in metric tubes of fixed
radius and in Newton-decrement neighborhoods of fixed radius
$0\leq\beta<1/2$, the least step count exceeds the unrestricted
optimum by a factor $\Theta(\sqrt{\log r})$.
\end{corollary}
\begin{proof}
Choose an absolute $\tau$ so that
$v(\tau+\log(4/5))\geq1/2$, with $v=H'$ from \eqref{eq:standard-profile}.
Equations \eqref{eq:radial-ordered-speed} and
\eqref{eq:radial-q-comparison} imply
\begin{equation}
 y_i'(s)\geq\kappa:=7/20\quad(s\geq a_i+\tau),
 \qquad 0\leq y_i(s)\leq H(s-a_i).
 \label{eq:radial-discrete-profile}
\end{equation}
The $U$ center at $T$ is accurate, so the radial center at
$T+\log(5/4)$ is accurate by \eqref{eq:radial-q-comparison}.
The Euclidean norm of its $\rho$ coordinates is $O(r\sqrt{\log r})$ by
\eqref{eq:dyadic-tail}. The straight segment from zero to this
endpoint in $\rho$ coordinates has $F$-length at most
$\sqrt{11/8}$ times its $U$-length, by
\eqref{eq:radial-metric-comparison}. Chord sampling of this segment
gives the explicit unrestricted sequence. Global metric domination and
the lower bound on $L_\opt(\eps_r)$ in \cref{thm:dyadic}
give the matching $\Omega(r\sqrt{\log r})$ distance.

For the lower bound, use thresholds $\sigma_i=a_i+\tau$,
$J=\lfloor r^{2/3}\rfloor$, cutoff $b=\sigma_{J+1}$, and
the potential
\[
 P(s)=\int_{\sigma_1}^{\min(\max(s,\sigma_1),b)}
       \sqrt{\#\{i:\sigma_i\leq u\}}\dd u.
\]
Its total range is $\Theta(r\log r)$, by the same harmonic
sum and bounded rounding error as in \cref{thm:tube}.
The step bound $R$ and the two tube conditions bound the $F$-distance
between consecutive labeled centers by $C_r=2\delta_r-\log(1-R)$.
By metric domination this also bounds the Euclidean distance
between their $y$ coordinates. All coordinates increase by
\eqref{eq:radial-ordered-speed}, so clipping labels does not
increase this distance. The first $j$ active coordinates advance
at rate at least $\kappa$, and successive thresholds
are separated by at least $\log2$. Thus the calculation
\eqref{eq:one-progress-jump}, with $v_0$ replaced by $\kappa$,
bounds a progress jump by $O(C_r\sqrt{1+C_r})$.

Initially, the direct prefix estimate in the proof of
\cref{thm:radial-upper} and \eqref{eq:radial-discrete-profile}
give
\[
 P(s_0)\leq\kappa^{-1}\int_{-\infty}^{s_0}\norm{y'(s)}_2\dd s
 \leq\kappa^{-1}\frac{25}{14}\Gamma_r\norm{y(s_0)}_2
 \leq\kappa^{-1}\frac{25}{14}\Gamma_r\delta_r.
\]
At termination $1-z_{N,1}\leq\eps_r$; scalar contraction and
the terminal tube imply
$y_1(s_N)\geq T-\log r-\delta_r$.
Together with $y_1(s)\leq H(s)\leq s_++1/\sqrt2$, this
forces $s_N>b$ for large $r$. Since
$T-b=\Theta(r^{2/3})$, the assumption $\delta_r=o(r^{2/3})$
suffices here.
The net progress is $\Theta(r\log r)$; summing the jump bound over
all steps proves \eqref{eq:growing-tube} without assuming monotone
labels.
Applying the same endpoint argument to the first accurate center
and integrating \eqref{eq:radial-discrete-profile} gives the
central-length lower bound; \cref{thm:radial-upper} gives its
upper bound. Finally apply \cref{lem:decrement-tube}.
\end{proof}

\begin{example}[A coupling singular at the vertices]
\label{ex:vertex-singular}
For $r\geq2$, consider $F_0=U-\log(r-\norm x^2)$. It has exact
parameter $r+1$, so it lies outside \eqref{eq:radial-barrier}.
Adding the parameters of the two terms gives the upper bound; along $x=(1-\tau)\mathbf1$
its restriction is $-(r+1)\log\tau+O(1)$ with the corresponding
first two derivative asymptotics, proving the reverse bound.
Its endpoint and equal-accuracy suprema are still at least
$\Gamma_{r-1}$. To prove this despite the singular coupling,
set $n=r-1$, $w_i=e^{Md_i}$ for $i\leq n$, and
$w_r=e^{-M^2}$, terminating at $s=0$. Stationarity gives
$0\leq x_r(s)\leq e^{-M^2}/2$ for $s\leq0$.
Thus $r-\norm x^2\geq1-x_r^2\geq1/2$ on this central arc
and on its coordinatewise $U$-geodesic from zero. The extra
gradient is bounded by $4$ and its Hessian by $(4+16r)I$ on
these paths. Apply \eqref{eq:coupled-prefix-speed} to the first
$n$ coordinates. Their prefix intervals give length
$M\Gamma_n^2-o_r(M)$; their endpoint metric coordinates are
$Md_i+O_r(1)$, while the final coordinate is $o(1)$.
The comparison geodesic has $F_0$-length at most
$M\Gamma_n+O_r(1)$, since its Euclidean total variation is at
most $r$. The ratio lower bound follows, and choosing the actual
terminal gap gives the lower bound on the equal-accuracy supremum
exactly as in
\cref{thm:facet-coupling}. Thus the lower bound does not need
\eqref{eq:facet-collar} on a whole facet collar: a coordinate with
negligible weight keeps the coupling regular along the central arc and
the comparison path.
\end{example}
